\section{Point-based evaluation and finite horizons}\label{sec:pbvi}

The earlier editions evaluated beliefs ``point-based'': for a horizon $k$, every word $w\in U_m^k$ produces the \emph{outcome mask} $A_k(w,L)=\{\theta:\theta\text{ survives }w\text{ for }k\text{ steps from }L\}$, and the $k$-step value of a belief is $\mathcal V^{(k)}_L(b)=\max_w b(A_k(w,L))$, a maximum of $b$-masses over the \emph{alpha-set} $\calA_k(L)=\{A_k(w,L)\}$ of distinct masks. This is exact for what it computes---$k$-step survival---and the set $\calA_k(L)$ is small: Table~\ref{tab:alpha} gives its size for $k\le4$ at $m=4$ ($17^4=83{,}521$ words, at most $857$ masks) and $m=5$ ($33^4\approx1.19\cdot10^6$ words, at most $15{,}929$ masks), reproducing the counts of the earlier scripts, and for $k\le3$ at $m=6$. What the earlier editions did not make explicit is the relation between $\mathcal V^{(k)}_L$ and the indefinite value $\mathcal V_L$ of Proposition~\ref{prop:antichain}.

\begin{proposition}[finite versus indefinite horizon]\label{prop:pbvi}
$\mathcal V^{(k)}_L(b)\ge\mathcal V^{(k+1)}_L(b)\ge\mathcal V_L(b)$ for all $k$, with $\mathcal V^{(k)}_L=\mathcal V_L$ for every belief as soon as every mask of $\calA_k(L)$ is in $\calS_m(L)$. For $m\le6$ this holds for $k\ge L+1$ (Corollary~\ref{cor:uniform}). At $m=4,6$ it fails for every $k\le L$ when $L\ge1$, because a far pair at distance $\delta_m+1$ survives exactly $L$ steps; at $m=5$ the longest survival of a non-survivable set is $2\lfloor(L-1)/2\rfloor$ steps (far pairs at distance $3$ and tetrahedra alike, Section~\ref{sec:criterion}), so the point-based value is exact precisely for $k>2\lfloor(L-1)/2\rfloor$: with $k=4$ this is $L\le4$ at $m=5$ against $L\le3$ at $m=4$.
\end{proposition}

Table~\ref{tab:alpha} also reports, for each horizon, how many masks are survivable sets: at $m=4$ and $L=15$ none of the $89$ two-step or $321$ three-step masks is, and $48$ of the $545$ four-step masks are, so a four-step evaluation at $z_0=2.5$ still contains $497$ masks that describe transient survival only. The supplement (S6) tabulates $\mathcal V^{(k)}_L$ for $k=1,\dots,4$ and $\mathcal V_L$ at probed beliefs---the uniform belief, five random beliefs drawn with a fixed seed, and three adversarial beliefs concentrated on a far pair, on an antipodal pair, and on a ball together with its antipode---at every level $L=0,\dots,15$. Three readings: for the uniform belief at $m=4$ the four-step value is $\tfr12$ at $L=15$ ($8$ of $16$ cells survive four steps) while the indefinite value is $\tfr18$ at every $L\ge1$; for the far-pair belief the four-step value is $1$ for every $L\ge4$ and the indefinite value $\tfr12$ at every level; at $m=5$ the uniform four-step value at $L=15$ is $\tfr12$ against $\tfr3{16}$. For $L\le3$ at $m=4$ and $L\le4$ at $m=5$ the finite and indefinite values agree for every probed belief, as Proposition~\ref{prop:pbvi} predicts, and they part at the next level. The point is structural: finite-horizon alpha-sets are exact evaluations of a different quantity, and the horizon at which they become exact for indefinite survival grows with the slack ($L+1$ at $m=4,6$, $2\lfloor(L-1)/2\rfloor+1$ at $m=5$); it is not $4$.

\begin{table}[t]
\centering\small
\caption{Size of the alpha-set $\calA_k(L)$ (distinct $k$-step outcome masks over all words), and in parentheses the number of those masks that are blind-survivable sets, for $k=4$ ($m=4,5$) and $k=3$ ($m=6$); the full $k=1,\dots,4$ tables are S6.1--S6.2.}
\label{tab:alpha}
\begin{tabular}{rlll}
\toprule
$L$ & $m=4$, $k=4$ & $m=5$, $k=4$ & $m=6$, $k=3$\\
\midrule
0 & 17 (16) & 945 (944) & 5{,}281 (5{,}280)\\
1 & 49 (48) & 945 (944) & 11{,}953 (11{,}952)\\
2 & 49 (48) & 945 (944) & --\\
3 & 49 (48) & 1{,}345 (1{,}344) & 33{,}649 (18{,}224)\\
4 & 209 (48) & 1{,}505 (1{,}504) & --\\
5 & 161 (48) & 5{,}265 (1{,}504) & 27{,}409 (13{,}584)\\
6 & 161 (48) & 5{,}265 (1{,}504) & --\\
7 & 185 (48) & 3{,}505 (1{,}344) & --\\
8 & 361 (48) & 6{,}065 (1{,}664) & --\\
9 & 297 (48) & 9{,}625 (1{,}664) & --\\
10 & 249 (48) & 7{,}705 (1{,}664) & 16{,}077 (2{,}144)\\
11 & 249 (48) & 6{,}937 (1{,}536) & --\\
12 & 857 (48) & 14{,}649 (1{,}696) & --\\
13 & 793 (48) & 15{,}929 (1{,}216) & --\\
14 & 729 (48) & 12{,}937 (1{,}104) & --\\
15 & 545 (48) & 11{,}145 (1{,}072) & 24{,}921 (0)\\
\bottomrule
\end{tabular}

\end{table}

\paragraph{What is superseded.} The auxiliary counts of the earlier editions---the number of two-cell supports met in a four-step evaluation, the reduction ``$17^4\to857$''---are properties of $\calA_4(L)$ and remain correct as such; the classification of Section~\ref{sec:classification} replaces them as the statement about survival. All $120$ pairs of the four-cube are now classified by distance ($32$ adjacent and survivable, $48$ at distance $2$, $32$ at distance $3$, $8$ antipodal), and the census tables count sets, not evaluations.
